%%%%%%%%%%%%%%%%%%%%%%%%%%%%%%%%%%%%%%%%%%
\section{Discussion}

\subsection{What the Coefficient-Space Network Does and Does Not Achieve}\label{sec:niche}

\rev{On the forward problem the answer to the question of whether a PINN can compete with a classical spectral solver is, in this study, no. On B1 and B2 Chebyshev collocation with an ordinary first-order time step is \N{b1.ratio.fh_K32_cheb80} and \N{b2.ratio.hard_cheb80} times more accurate than the revised Spectral-PINN and several orders of magnitude cheaper; on B1 a model that ignores diffusion altogether is within a factor \N{b1.ratio.ode_fh_K32} of the network; and on B3 neither PINN learns the dynamics at the budget used here, while the classical solver does. This agrees with the established outcome for forward problems~\cite{ref_Grossmann2024,ref_McGreivy2024}, and we do not offer the method as a competitor on this task. What the experiments do show is narrower. (i)~Inside the family of PINNs, putting the discretization into the output space removes spatial automatic differentiation and gives, on B1, a gain over the space--time PINN of \N{b1.ratio.van_fhK32} in accuracy (\N{b1.ratio.lb} when both are refined by L-BFGS) and of about \N{b1.speed.van_fhK32} in training time (Adam phase); on B2 the gain is \N{b2.ratio.van_hard} with Adam alone and disappears after L-BFGS (one seed each). (ii)~Structure that is a linear constraint on the coefficients can be imposed exactly: the Neumann basis~\eqref{eq:neumann-basis} does so at no cost in accuracy and removed a failure that a penalty produced in both of its scalings (Section~\ref{sec:results-k}), and the reflection-symmetric Fourier basis keeps the symmetry violation at round-off (Section~\ref{sec:2d-results}). (iii)~The energy of the learned trajectory is controlled a posteriori by quantities computable to round-off (Theorem~\ref{thm:balance}), a certificate that a standard PINN lacks. Such a network could be useful where a classical solver is awkward: a continuous-in-time representation differentiable in its parameters, coupling to data when part of the physics is unknown, and parametric families~\cite{ref_Du2024}. We tested none of these, and our inverse experiments give no support for data assimilation in the simplest setting: with a known model structure, classical least squares is at least as accurate as the best neural estimator on noise-free data (better than $10^{-3}\%$) and much more accurate on noisy data. The contribution of this article is a controlled test bed, a set of diagnostics and design principles, and a list of pitfalls that we encountered (and in part fell into), not a new solver.}

\subsection{Pitfalls Found in the First Version}

\rev{Three conclusions of the first version did not survive a systematic check, and each rests on a mistake that is easy to make. (i)~\emph{Large loss of accuracy with growing $K$} was caused by a boundary penalty whose effective weight decayed like $K^{-4}$. Rescaling the penalty did not cure it (the loss becomes ill-conditioned); the exact Neumann basis did. The milder growth of the error with $K$ that remains at a fixed budget is a matter of optimization: twice the epochs reduce the error of seed~42 from \N{kk.hard_cgl64_s42} to \N{kk.kh_K64_x2.l2} at $K=64$ but only from \N{kk.hard128_s42} to \N{kk.kh_K128_x2.l2} at $K=128$. The wall slope of the trained network is a diagnostic that would have caught the penalty failure. (ii)~\emph{A systematic $38\%$ bias of the inverse estimate} reflects the start and the step size of the optimization, not the data; a result that scales with an arbitrary starting value signals that the optimizer, not the estimator, is being measured. (iii)~\emph{Zero energy violations as an effect of the energy penalty.} The unpenalized network has none either; the first version reported the count only for the penalized one, so the missing control was the comparison of the two counts. In addition, B1 is reaction-dominated and almost independent of $\varepsilon$ for $\varepsilon\le0.01$, so that the small-$\varepsilon$ experiments tested essentially one function, and the initial condition did not satisfy the boundary condition. These flaws were found through a reviewer's remark, an audit of every number of the first version against its data (22 of 77 numerical statements were wrong and 7 could not be verified) and a certified reference solution; we recommend the same checks to anyone benchmarking PINNs.}

\subsection{Energy Decay and Admissibility}\label{sec:energy-discussion}

\rev{Monotone energy is a necessary but weak property. The energy of every one of the \N{en.n_all_spec} Spectral-PINNs trained for the forward problem decreased monotonically, with or without the penalty, including networks whose $L^2$ error is of order one and which violate $|u|\le1$ by up to \N{mbp.excess_all}; accurate networks overshoot by at most a few $10^{-3}$ (revised Spectral-PINNs of B1 outside the poor-trajectory regimes by up to \N{mbp.excess}), and the vanilla PINN, which has no structural guarantee, also has a decreasing energy on B1 and B2 (\N{en.van_viol} increases in the \N{mbp.n_van} runs on B1, on the grid of spacing $0.02$). Only on B3 does the vanilla PINN show energy increases (\N{en.b3van_inc} of \N{en.b3van_int} intervals), in a run whose error is of order one anyway. Theorem~\ref{thm:balance} explains half of this: whatever the network does, its energy can rise only through the PDE residual and the boundary defect, so a network that is good in the PDE sense is nearly dissipative and the penalty has nothing to add. The converse does not hold: a decreasing energy says nothing about accuracy, and the maximum bound has to be checked separately. The theorem does not say that a trained network is stable or convergent; a convergence bound in terms of the optimization error would need, in addition, a stability estimate for the PDE (as in~\cite{ref_Mishra2023}), which we do not provide.}

\subsection{Limitations}

\rev{The study is limited to a one-dimensional equation and one two-dimensional symmetry test, to CPU training of small networks, and to five seeds for the main configurations (one to three for the others; see the $n$ column of the tables). On B3 both PINNs fail with the plain training used here (at $t=0.2$ their errors are already about $5\times10^{-2}$ for the Spectral-PINNs and $10^{-1}$ for the vanilla PINN, i.e.\ the optimization converges to a wrong trajectory); remedies such as causal training, curriculum in time or time marching~\cite{Wang_causal,Mattey} were not tried. The vanilla baseline is plain, without causal training, sequential training, adaptive sampling~\cite{Wight} or Fourier features~\cite{Liu_FourierBias2025}, to which the literature attributes large gains; the Spectral-PINN is also plain. Wall-clock times of the training runs come from a shared server and are indicative only. The inverse experiments use synthetic data from the same equation and a single network initialization per configuration. The energy penalty was studied only for a gradient flow with a monotone energy.}
